\subsection{Weighted transfer for a nonunit singleton}\label{proof:Q5-E008-singleton}
\paragraph{The square-root multiplier.}
Write
\[
A=J+K,\qquad T=A^*=J-K,\qquad g=K\one,
\qquad p=1+g,\quad q=1-g,
\]
\[
f_n=A^n\one,\qquad\bar f_n=T^n\one,
\qquad\delta_j=\lVert K^j\one\rVert_2^2,\quad d=\delta_1.
\]
Let \(M_v\) denote multiplication by \(v\), and set
\(I_q(Y)=\E[q^{-1/3}|Y|^{8/3}]\). The square-root estimate
\eqref{q3:eq:half}, including every \(a\ge2\), is
\[
\lVert M_{\sqrt{f_a}}T\rVert_{2\to2}\le1.
\]
Replace \(K\) by \(-K\) in the entire construction to obtain
\(\lVert M_{\sqrt{\bar f_a}}A\rVert_{2\to2}\le1\).
Taking the ordinary adjoint on the original measure gives
\[
\boxed{\lVert T M_{\sqrt{\bar f_a}}\rVert_{2\to2}\le1,
\qquad a\ge2.}
\tag{E8.10.1}
\]
The displayed multiplier is a square root.

\paragraph{Full reversal of the actual reference input.}
Use the complete outward-\(f_a\) reference directions in
Sections~\ref{proof:R2-C032} and \ref{proof:R2-C034}. They give the
reference fourth moment for every nonreducible terminal and every
original inner length \(L\ge1\), with bound \(1-d/3000\) for \(a\ge3\)
and bound one for \(a=2\). Their actual direction tables are used with
all their exceptions. In particular, the unit-\((1,\allowbreak 2,\allowbreak 3)\) reference
exception is at \(L=1\); the ordinary unit-anchor eighth-moment exception at \(L=2\)
has a different role.

Let \(F_{\rm ref}\) denote the actual terminal selected by those
reference directions. Reverse the entire reference construction:
the actual reference arm, terminal arms, all \(L\) connector edges and
both decoration products. A superscript minus denotes this complete
reversal. Put
\[
\begin{gathered}
X_0=A^L(E_t^-F_{\rm ref}^-),\qquad X=E_r^-X_0,\\
Q_a(d)=
\begin{cases}
1,&a=2,\\
1-d/3000,&a\ge3.
\end{cases}
\end{gathered}
\]
Since \(0<E_r^-\le1\), the fully reversed reference bound gives
\[
\E[\bar f_aX^4]
\le\E[E_r^-\bar f_aX_0^4]\le Q_a(d).
\tag{E8.10.2}
\]
Here each root equal pair occurs once, and every terminal pair remains
inside its actual \(A^L\) connector. The sole nonunit residual root arm
points inward and gives \(\bar f_a\). The actual final component is
\[
U=T^M(\bar f_aX).
\tag{E8.10.3}
\]
Thus the inner connector points from the root toward its terminal, the
root arm toward the root, and the outer connector from the root toward
the center. These directions are fixed by the original structure.
Every arm in \(F_{\rm ref}^-\) is reversed together with the reference
and connector: for example, \(qf_b\) becomes \(p\bar f_b\), and
\(f_3\bar f_b\) becomes \(\bar f_3f_b\). The output weight \(q\)
belongs to the actual outer \(T\)-path.

\paragraph{The weighted fourth-moment bridge.}
For any real \(X\), \(a\ge2\) and \(M\ge1\), consider the positive
operator \(P_0=T^MM_{\bar f_a}\). Its row is
\(P_0\one=\bar f_{M+a}\). Row Cauchy--Schwarz gives
\[
|P_0X|^2\le\bar f_{M+a}\,T^M(\bar f_a|X|^2).
\]
Apply H\"older with exponents \(3\) and \(3/2\), then cube, to obtain
\[
I_q(P_0X)^3
\le\E[\bar f_{M+a}^4/q]
 \left\{\E\left[(T^M(\bar f_a|X|^2))^2\right]\right\}^2.
\]
In \(L^2\), factor the vector as
\[
T^M(\bar f_a|X|^2)
=[T^MM_{\sqrt{\bar f_a}}]\,[\sqrt{\bar f_a}|X|^2].
\]
Consequently the exact weighted bridge is
\[
\boxed{
\begin{aligned}
I_q(T^M(\bar f_aX))^3
\le{}&\E[\bar f_{M+a}^4/q]\,
 \lVert T^MM_{\sqrt{\bar f_a}}\rVert_{2\to2}^4\\
&\hspace{2em}\cdot(\E[\bar f_a|X|^4])^2.
\end{aligned}}
\tag{E8.10.4}
\]
Both the norm and input weight follow from this factorization. The
factorization uses the one \(\bar f_a\) arm as two square-root multipliers.

To keep every operator coefficient explicit, write
\[
\begin{gathered}
N_0=1,\qquad N_1=1+(32/19)d,\qquad
N_2=(1+(1024/855)\delta_2)^2,\\
N_n=(1-d/8)^{\lfloor n/3\rfloor}\qquad(n\ge3).
\end{gathered}
\]
Section~\ref{proof:Q5-C064} proves \(\lVert T^n\rVert\le N_n\).
These \(N_n\) bound the norm itself. The bases are the separate estimates
for powers three, four and five; writing lengths as \(3m,\allowbreak 3m+1,\allowbreak 3m+2\)
uses respectively \(m\) cubes, one fourth power with \(m-1\) cubes, or
one fifth power with \(m-1\) cubes.  Use (E8.10.1) in the actual order
\[
T^MM_{\sqrt{\bar f_a}}
=T^{M-1}(TM_{\sqrt{\bar f_a}}),\qquad
\lVert T^MM_{\sqrt{\bar f_a}}\rVert\le N_{M-1}.
\]
The bare inverse-row estimates in Sections~\ref{proof:Q5-C064}
and \ref{proof:Q5-C100} give
\(\E[\bar f_n^4/q]\le1-d/2\) for every \(n\ge3\).
Their \(n=3\) bound is the separate \(1-23d/36\) estimate; the
\(n\ge4\) row uses the complete longer-path proof. Combining these
with (E8.10.1), (E8.10.2) and (E8.10.4) yields
\[
\boxed{I_q(U)^3\le(1-d/2)N_{M-1}^4Q_a(d)^2.}
\tag{E8.10.5}
\]
For \(M=1\), the multiplier norm is at most one. For \(M\ge4\), the
remaining exponent \(M-1\ge3\) is covered by the separate power bases.
Thus
\[
I_q(U)^3\le1-d/2\le e^{-d/2}\qquad(M=1),
\tag{E8.10.6}
\]
\[
I_q(U)^3\le
e^{-[1/2+\lfloor(M-1)/3\rfloor/2]d}\qquad(M\ge4).
\tag{E8.10.7}
\]
For \(a\ge3\), the additional factor \(Q_a(d)^2\) may be retained.
The displayed common coefficients cover every nonunit singleton,
nonreducible terminal and inner length \(L\) in these two outer ranges.
At \(a=2,\allowbreak M=1\), the actual row is \(\bar f_3\) and uses its separate
bound; every other relevant row has length at least four.
At \(d=0\), the operator proof still applies, and actual leaf integration
gives equality.
